\documentclass[10pt,letterpaper]{article}
\usepackage[margin=0.55in]{geometry}
\usepackage[T1]{fontenc}
\usepackage[utf8]{inputenc}
\usepackage{lmodern}
\usepackage{xcolor}
\usepackage[hidelinks]{hyperref}
\usepackage{tabularx}
\usepackage{array}
\pagestyle{empty}
\setlength{\parindent}{0pt}
\setlength{\parskip}{0pt}
\definecolor{accent}{HTML}{174A7C}
\definecolor{textgray}{HTML}{333333}
\newcommand{\resumeHeader}[4]{%
  {\Huge\bfseries #1}\\[-1pt]
  {\large\bfseries\color{accent}#2}\\[3pt]
  {#3 \; | \; #4}\\[-2pt]
}
\newcommand{\role}[4]{%
  \textbf{#1} \hfill #2\\
  \textit{#3} \hfill \textit{#4}\vspace{2pt}
}
\newcommand{\resumelist}{\begin{itemize}\setlength{\itemsep}{1.2pt}\setlength{\parskip}{0pt}\setlength{\parsep}{0pt}}
\newcommand{\resumelistend}{\end{itemize}}
\renewcommand\labelitemi{\textcolor{accent}{\small$\bullet$}}
\begin{document}
\color{textgray}

\resumeHeader{YASH DESHPANDE}{Principal Machine Learning Engineer | AI Systems | Multimodal AI | Computer Vision}{Boston, MA | +1 917 362 9676 | \href{mailto:yashrajas@gmail.com}{yashrajas@gmail.com}}{\href{https://www.linkedin.com/in/yash--deshpande/}{linkedin.com/in/yash--deshpande} | \href{https://github.com/yashd94}{github.com/yashd94}}

\vspace{4pt}\textcolor{accent}{\large\bfseries SUMMARY}\vspace{2pt}\hrule\vspace{4pt}
Principal-level Machine Learning Engineer and Applied AI Scientist with 8+ years of experience building scalable ML platforms, multimodal AI systems, computer vision models, NLP solutions, and cloud-native machine learning infrastructure. Experienced leading end-to-end ML initiatives across research and production settings, including large-scale image processing, representation learning, Bayesian modeling, high-dimensional data integration, reproducible experimentation, and GPU-accelerated inference workflows. Strong technical leader who translates ambiguous scientific and product objectives into robust ML architecture, production-grade pipelines, and reusable engineering systems.

\vspace{6pt}\textcolor{accent}{\large\bfseries CORE STRENGTHS}\vspace{2pt}\hrule\vspace{4pt}
\begin{tabularx}{\textwidth}{>{\bfseries\color{accent}}p{1.15in} X}
AI/ML & Deep learning, computer vision, foundation models, multimodal learning, representation learning, NLP, graph neural networks, Bayesian ML, statistical learning, predictive modeling \\
ML Systems & Production ML systems, distributed data processing, GPU inference, feature pipelines, experiment tracking, model monitoring, reproducible ML workflows, performance optimization \\
Engineering & Python, PyTorch, TensorFlow, scikit-learn, pandas, NumPy, SQL, R, Shell, Docker, Git, MLflow, Weights \& Biases, Nextflow, CI/CD \\
Cloud/Data & AWS, Google Cloud Storage, large-scale imaging data, Parquet, Zarr, DuckDB, genomics, transcriptomics, histopathology, clinical text \\
Leadership & Technical strategy, platform architecture, cross-functional execution, stakeholder alignment, engineering standards, mentoring, roadmap definition
\end{tabularx}

\vspace{6pt}\textcolor{accent}{\large\bfseries PROFESSIONAL EXPERIENCE}\vspace{2pt}\hrule\vspace{4pt}
\role{Merck \& Co.}{Dec 2023 - Present}{Senior Scientist, Machine Learning}{Boston, MA}
\resumelist
  \item Lead machine learning initiatives spanning multimodal AI, single-cell foundation modeling, and computer vision for large-scale biomedical data analysis.
  \item Architect ML systems for genetic perturbation response prediction using single-cell RNA-seq datasets, linking learned representations to target discovery and prioritization workflows.
  \item Design scalable whole-slide image processing and embedding pipelines for H\&E histopathology, including patch extraction, tissue masking, coordinate stores, GPU inference, and downstream analytics outputs.
  \item Build and optimize high-throughput data pipelines using PyTorch, Parquet, Zarr, GCS, local NVMe caching, batched transfer workflows, and resumable processing patterns for large imaging corpora.
  \item Develop computer vision models for cell segmentation and tissue characterization, enabling quantitative imaging biomarkers for translational and companion diagnostic applications.
  \item Establish reproducible ML development patterns, experiment tracking, data validation, logging, and performance diagnostics for long-running, large-scale training and inference workflows.
  \item Partner with translational scientists, computational biologists, ML engineers, and research stakeholders to define technical strategy and convert exploratory modeling efforts into reusable AI systems.
\resumelistend

\role{Prometheus Biosciences, a wholly owned subsidiary of Merck \& Co.}{Jun 2022 - Dec 2023}{Machine Learning Engineer}{San Diego, CA}
\resumelist
  \item Led development of \textbf{Minerva}, a multimodal companion diagnostics discovery platform within Prometheus360, enabling versioned, reproducible, and extensible model development across heterogeneous omics datasets.
  \item Architected scalable ML workflows for feature generation, whole-genome scanning, QTL analysis, evidence meta-analysis, and biomarker prioritization using Nextflow and cloud-native processing patterns.
  \item Built Bayesian classification and genotype clustering workflows for predictive modeling downstream of SNP discovery and companion diagnostic feature selection.
  \item Created reusable infrastructure for model development, provenance tracking, reproducibility, and collaborative experimentation across scientific and engineering teams.
  \item Acted as a technical bridge between ML engineering, immunology, computational biology, radiology, and biomarker strategy teams to align platform capabilities with high-value translational use cases.
\resumelistend

\role{Fresenius Medical Care North America}{Oct 2020 - May 2022}{Machine Learning Engineer}{Remote / Boston, MA}
\resumelist
  \item Built cloud-based ML platform capabilities for an internal AI workbench on AWS, enabling scalable model monitoring and faster iteration across cross-functional analytics teams.
  \item Developed reusable ML infrastructure and tooling to accelerate adoption of cloud-native experimentation, deployment, and monitoring workflows.
  \item Built a Python-based NLP model for unstructured clinical text that identified 40+ previously unreported annual abdominal pain events, improving downstream clinical signal detection.
  \item Improved ML engineering velocity by standardizing infrastructure patterns, data access workflows, and operational monitoring for applied AI initiatives.
\resumelistend

\role{ZS Associates}{Jul 2016 - May 2018}{Data Analyst, Business Technology Group}{Pune, India}
\resumelist
  \item Built and maintained production ETL pipelines using Shell, PostgreSQL, AWS Redshift, and workflow automation for client-facing analytics platforms.
  \item Reduced front-end processing time by 70\% on a business rules platform through performance-focused workflow redesign and automation.
  \item Supported recurring analytics and decision-support systems through scalable data engineering, query optimization, and operational process improvements.
\resumelistend

\vspace{6pt}\textcolor{accent}{\large\bfseries SELECTED PRINCIPAL ML ENGINEERING HIGHLIGHTS}\vspace{2pt}\hrule\vspace{4pt}
\resumelist
  \item Built end-to-end AI systems across data ingestion, feature engineering, model development, training/inference orchestration, experiment tracking, and downstream analytics integration.
  \item Designed large-scale computer vision pipelines for whole-slide histopathology data, including patch coordinate extraction, tissue quality metrics, embedding generation, local/remote storage synchronization, and cache eviction.
  \item Delivered multimodal ML platform capabilities supporting molecular, imaging, genotype, and clinical data integration for high-stakes decision workflows.
  \item Developed NLP, Bayesian ML, graph neural network, and representation learning solutions across clinical text, imaging, molecular, and multimodal datasets.
  \item Established engineering practices for resumability, observability, validation, and performance diagnostics in long-running ML pipelines.
\resumelistend

\vspace{6pt}\textcolor{accent}{\large\bfseries EDUCATION}\vspace{2pt}\hrule\vspace{4pt}
\textbf{New York University} \hfill Sep 2018 - May 2020\\
Master of Science in Data Science\\[2pt]
\textbf{University of Pune} \hfill Aug 2012 - May 2016\\
Bachelor of Engineering in Electronics \& Telecommunication Engineering

\vspace{6pt}\textcolor{accent}{\large\bfseries SELECTED PUBLICATIONS}\vspace{2pt}\hrule\vspace{4pt}
\resumelist
  \item Gerardo A. Flores et al. \textit{Pandora's Regret: A Proper Scoring Rule for Evaluating Sequential Search}. arXiv:2605.01936, 2026.
  \item Ruiwen Ding et al. \textit{Predicting Ulcer in H\&E Images of Inflammatory Bowel Disease Using Domain-Knowledge-Driven Graph Neural Network}. arXiv:2504.09430, 2025.
  \item Houda Alberts et al. \textit{VisualSem: A High-Quality Knowledge Graph for Vision and Language}. arXiv:2008.09150, 2020.
  \item Ningyuan Huang et al. \textit{Endowing Language Models with Multimodal Knowledge Graph Representations}. arXiv:2206.13163, 2022.
\resumelistend

\vspace{6pt}\textcolor{accent}{\large\bfseries LANGUAGES}\vspace{2pt}\hrule\vspace{4pt}
English, Marathi, Hindi; basic Japanese

\end{document}
